\documentclass{article}
\usepackage{amsmath}
\usepackage{amsfonts}


\newcommand{\0}{\mathbf{0}}
\newcommand{\Bdiag}{\text{Bdiag}}
\newcommand{\pen}[2]{\lambda#1+#2}
\newcommand{\penEA}{\pen{E}{A}}
\newcommand{\C}{\mathbb{C}}
\newcommand{\N}{\mathbb{N}}
\renewcommand{\vec}{\mathbf}
\newcommand{\inv}{^{-1}}

\begin{document}
Two matrix pencils $\lambda A + B$, $\lambda C + D$ are similar if there exists matrix $P, Q$ s.t. $\lambda A + B = P (\lambda C + D) Q$.
% Not relevent in this case, but do P,Q need to be invertible in general? Do the pencils need to be regular?
Note that similar pencils have similar determinant.
% If two polynominals differ by a factor, is that refered to as similar? Or is there another name?

Given $n$ by $n$ matrices $E$ and $A$, such that the matrix pencil $\lambda E + A$ is regular ($|\lambda E+A| \neq \0$), our goal is to find $P,Q \in M_n$ s.t.
\begin{equation*}
    P (\lambda E + A) Q = \Bdiag(\lambda I - J, \lambda N - I)
\end{equation*}

Let $n_J$ be the degree of the determinant of the pencil. We claim the matrix pencil $\penEA$ is similar to a matrix of the form
\begin{equation}
    \label{equ:half_upper_triangular_general}
    \left[\begin{array}{c|c}
        \begin{matrix}
            \lambda-\lambda_1&\cdots&&\\
            &\lambda-\lambda_2&\cdots&\\
            &&\ddots&\cdots\\
            &&&\lambda-\lambda_{n_J}\\
        \end{matrix}&\cdots\\
        \hline\\
        &\cdots\\
    \end{array}\right]
\end{equation}

We will show this by induction on $n_J$.
If $n_J$ is $0$, we are done.
Otherwise, suppose it holds true for $n_J-1 \in \N$.
Since the pencil has determinant of degree at least one, there exists at least one root for the equation $|\penEA|=0$.
Let $\lambda_1$ be one such root, we can find vector $\vec{v_1} \in \C^n \neq \0$ such that $(\lambda_1E+A)\vec{v_1}=\0$.
We can complete the set $\{\vec{v_1}\}$ to form a basis $\{\vec{v_1},\vec{v_2},\dots,\vec{v_n}\}$ for $\C^n$.
Let $V=\left[\begin{matrix}\vec{v_1}&\vec{v_2}&\dots&\vec{v_n}\end{matrix}\right]$.
Partition $V=\left[\begin{matrix}\vec{v_1}&\tilde{V}\end{matrix}\right]$.
Similarly, let $\vec{u_1}=E\vec{v_1}$.
If $\vec{u_1}=\0$, $A\vec{v_1}=\0-\lambda_1E\vec{v_1}=\0$ and $(\penEA)\vec{v_1}=\0$, contradicting the requirement for the pencil to be regular.
Hence, we cal also complete $\{\vec{u_1}\}$ to a basis of $C^n$ $\{\vec{u_1},\vec{u_2},\dots,\vec{u_n}\}$.
Let $U=\left[\begin{matrix}\vec{u_1}&\vec{u_2}&\dots&\vec{u_n}\end{matrix}\right]$. Partition $U=\left[\begin{matrix}\vec{u_1}&\tilde{U}\end{matrix}\right]$.
Note that \[A\vec{v_1}=-\lambda_1E\vec{v_1}=\lambda_1\vec{u_1}\].
We have
\begin{align*}
    \frac{1}{\lVert\vec{u_1}\rVert^2}U^*(\lambda E + A)V
    &=\frac{1}{\lVert\vec{u_1}\rVert^2}\begin{bmatrix}\vec{u_1}^*\\\tilde{U}^*\end{bmatrix}(\lambda E+A)\begin{bmatrix}\vec{v_1}&\tilde{V}\end{bmatrix}\\
    &=\frac{1}{\lVert\vec{u_1}\rVert^2}\begin{bmatrix}\vec{u_1}^*(\lambda E+A)\vec{v_1}&\vec{u_1}^*(\lambda E+A)\tilde{V}\\\tilde{U}^*(\lambda E+A)\vec{v_1}&\tilde{U}^*(\lambda E+A)\tilde{V}\end{bmatrix}\\
    &=\frac{1}{\lVert\vec{u_1}\rVert^2}\begin{bmatrix}\vec{u_1}^*(\lambda-\lambda_1)(E\vec{v_1})&\vec{u_1}^*(\lambda E+A)\tilde{V}\\\tilde{U}^*(\lambda-\lambda_1)(E\vec{v_1})&\tilde{U}^*(\lambda E+A)\tilde{V}\end{bmatrix}\\
    &=\begin{bmatrix}(\lambda-\lambda_1)&\frac{1}{\lVert\vec{u_1}\rVert^2}\vec{u_1}^*(\lambda E+A)\tilde{V}\\0&\frac{1}{\lVert\vec{u_1}\rVert^2}\tilde{U}^*(\lambda E+A)\tilde{V}\end{bmatrix}\\
\end{align*}

Let $\lambda E'+A' = \frac{1}{\lVert\vec{u_1}\rVert^2}\tilde{U}^*(\lambda E+A)\tilde{V}$.
Note that $|\lambda E+A|=(\lambda-\lambda_1)|\lambda E'+A'|$, so the degree of $|\lambda E'+A'|$ is $n_J-1$.
Thus, there exists $P,Q \in M_{n-1}$ s.t. $P(\lambda E'+A')Q$ is in the form of (\ref{equ:half_upper_triangular_general}).
Then,
\begin{align*}
    \begin{bmatrix}1&0\\0&P\end{bmatrix}\frac{1}{\lVert\vec{u_1}\rVert^2}U^*(\lambda E + A)V\begin{bmatrix}1&0\\0&Q\end{bmatrix}
    &=\begin{bmatrix}1&0\\0&P\end{bmatrix}\begin{bmatrix}(\lambda-\lambda_1)&\frac{1}{\lVert\vec{u_1}\rVert^2}\vec{u_1}^*(\lambda E+A)\tilde{V}\\0&(\lambda E'+A')\end{bmatrix}\begin{bmatrix}1&0\\0&Q\end{bmatrix}\\
    &=\begin{bmatrix}(\lambda-\lambda_1)&\frac{1}{\lVert\vec{u_1}\rVert^2}\vec{u_1}^*(\lambda E+A)\tilde{V}\\0&P(\lambda E'+A')\end{bmatrix}\begin{bmatrix}1&0\\0&Q\end{bmatrix}\\
    &=\begin{bmatrix}(\lambda-\lambda_1)&\frac{1}{\lVert\vec{u_1}\rVert^2}\vec{u_1}^*(\lambda E+A)\tilde{V}Q\\0&P(\lambda E'+A')Q\end{bmatrix}\\
\end{align*}
is also in the form of (\ref{equ:half_upper_triangular_general}).

We have shown that $\lambda E + A$ is similar to
\begin{equation}
    \label{equ:half_upper_triangular}
    \begin{bmatrix}
        T_1\lambda+T_2&M_1\lambda+M_2\\
        0&N_1\lambda+N_2\\
    \end{bmatrix}
\end{equation}
where $T_1, T_2$ are upper triangular and all elements on the diagonal of $T_1$ are 1s.
We also have $|N_1\lambda+N_2|=C$ for some $C\in\C\setminus\{0\}$.
We claim $N_2$ is invertible. Elegant proof pending.

To show this, we can left-multiply $N_1\lambda+N_2$ by a series of elementary row matrices $R$ to convert $N_2$ to row echleon form. If $N_2$ is singular, the last row will only include elements from $RN_1\lambda$. Thus, $\lambda$ divides $|N_1\lambda+N_2|=C$, which means $C=0$.

Since $N_2$ and $T_1$ are invertible, we can multiply (\ref{equ:half_upper_triangular_general}) by $\begin{bmatrix}
    T_1^{-1}&0\\0&N_2^{-1}
\end{bmatrix}$
Note that $T_1^{-1}$ is also upper triangular, so $T_2T_1^{-1}$ is still upper triangular.
Thus, we end up with
\begin{equation}
    \label{equ:half_upper_triangular_diag}
    \begin{bmatrix}
        I\lambda+T_3&M_1\lambda+M_2\\
        0&N_3\lambda+I\\
    \end{bmatrix}
\end{equation}

We can convert $T_3$ and $N_3$ to jordan form. Let $J=X_1\inv T_3X_1$ and $N=X_2\inv N_3X_2$.
\begin{align*}
    \begin{bmatrix}X_1\inv&0\\0&X_2\inv\\\end{bmatrix}\begin{bmatrix}I\lambda+T_3&M_1\lambda+M_2\\0&N_3\lambda+I\\\end{bmatrix}\begin{bmatrix}X_1&0\\0&X_2\\\end{bmatrix}
    &=\begin{bmatrix}X_1\inv(I\lambda+T_3)&X_1\inv(M_1\lambda+M_2)\\0&X_2\inv(N_3\lambda+I)\\\end{bmatrix}\begin{bmatrix}X_1&0\\0&X_2\\\end{bmatrix}\\
    &=\begin{bmatrix}X_1\inv(I\lambda+T_3)X_1&X_1\inv(M_1\lambda+M_2)X_2\\0&X_2\inv(N_3\lambda+I)X_2\\\end{bmatrix}\\
    &=\begin{bmatrix}I\lambda+J&M_3\lambda+M_4\\0&N\lambda+I\\\end{bmatrix}
\end{align*}

$N$ is nil potent, since otherwise the degree of $|N\lambda+I|$ will be non-zero.

\begin{align*}
    \begin{bmatrix}I&-M_4\\0&I\end{bmatrix}\begin{bmatrix}I\lambda+J&M_3\lambda+M_4\\0&N\lambda+I\\\end{bmatrix}
    &=\begin{bmatrix}I\lambda+J&(M_3-M_4N)\lambda\\0&N\lambda+I\\\end{bmatrix}\\
    &=\begin{bmatrix}I\lambda+J&M\lambda\\0&N\lambda+I\\\end{bmatrix}\\
\end{align*}

Since $N$ is nil potent, there exist minimum $m$ s.t. $N^m=0$.
Let
\begin{equation*}
    X = \sum_{i=0}^{m-1}J^iMN^i = M + JMN + J^2MN^2 + J^3MN^3 + \dots + J^{m-1}MN^{m-1}
\end{equation*} \footnote{We define $X^0=I$ for all $X$, not just non-singular one}
We have
\begin{align*}
    JXN
    &= J(\sum_{i=0}^{m-1}J^iMN^i)N\\
    &= \sum_{i=1}^{m}J^iMN^i\\
    &= \sum_{i=0}^{m-1}J^iMN^i + J^mMN^m - J^0MN^0\\
    &= X - M\\
\end{align*}


\begin{align*}
    \begin{bmatrix}I&JX\\0&I\end{bmatrix}\begin{bmatrix}I\lambda+J&M\lambda\\0&N\lambda+I\\\end{bmatrix}\begin{bmatrix}I&-X\\0&I\end{bmatrix}
    &=\begin{bmatrix}I\lambda+J&(M+JXN)\lambda+JX\\0&N\lambda+I\\\end{bmatrix}\begin{bmatrix}I&-X\\0&I\end{bmatrix}\\
    &=\begin{bmatrix}I\lambda+J&(M+JXN-X)\lambda+JX-JX\\0&N\lambda+I\\\end{bmatrix}\\
    &=\begin{bmatrix}I\lambda+J&(M+X-M+X)\lambda\\0&N\lambda+I\\\end{bmatrix}\\
    &=\begin{bmatrix}I\lambda+J&0\\0&N\lambda+I\\\end{bmatrix}\\
\end{align*}

\end{document}

